\clearpage
\begin{summarybox}

\textbf{五步认知链（先理解，再速算）：}
\begin{center}
\small 垂径定理 $\to$ 圆的中点弦斜率 $\to$ 椭圆的尺度变化\\[2pt]
$\to$ 双曲线的符号变化 $\to$ 点差法统一证明
\end{center}
\textbf{中点弦的最本质分工：}
\[
\boxed{\underbrace{\text{两端点方程相减}}_{\text{方向：往哪里走}}}
\quad+\quad
\boxed{\underbrace{\text{两端点方程相加}}_{\text{位移：走多远、能否到达}}}.
\]
\textbf{圆：}相减 $\Rightarrow OM\perp PQ$；相加 $\Rightarrow OM^2+MP^2=R^2$（勾股定理）。
\textbf{双曲线：}相减 $\Rightarrow x_0u/a^2=y_0v/b^2$；
相加 $\Rightarrow S+(u^2/a^2-v^2/b^2)=1$（\textbf{带符号}类比，非普通长度平方）。
\[
\boxed{u=a^2y_0t,\ v=b^2x_0t}
\quad\Longrightarrow\quad
\boxed{S-a^2b^2St^2=1}.
\]
$M\ne O$ 且 $S\ne0$ 时，$t^2=(S-1)/(a^2b^2S)$。
\textbf{$0<S<1$ 时从 $S$ 出发沿候选方向只会下降，走不到双曲线要求的 $1$；}
$S<0$ 时可上升到 $1$；$S>1$ 时可下降到 $1$。
原点是方向不唯一的例外。
\[
\boxed{\text{非退化实中点弦存在}\iff S<0\ \text{或}\ S>1\ \text{或}\ M=O}.
\]
\[
\begin{array}{c|c}
\text{曲线（均须存在非退化弦，且 }y_0\ne0\text{）}&\text{中点弦斜率}\\\hline
x^2+y^2=R^2 & -x_0/y_0\\[2pt]
 x^2/a^2+y^2/b^2=1 & -b^2x_0/(a^2y_0)\\[2pt]
 x^2/a^2-y^2/b^2=1 & +b^2x_0/(a^2y_0)
\end{array}
\]
\textbf{统一的最小推理：}
\[
\frac{x^2}{a^2}+\varepsilon\frac{y^2}{b^2}=1
\quad\xrightarrow{\text{端点平方差}}\quad
\boxed{\frac{x_0}{a^2}+\varepsilon\frac{y_0}{b^2}k=0}
\quad\Longrightarrow\quad
\boxed{k=-\varepsilon\frac{b^2x_0}{a^2y_0}}.
\]
椭圆 $\varepsilon=+1$，双曲线 $\varepsilon=-1$，圆取 $a=b=R,\varepsilon=+1$。
\textbf{注意：}统一公式求\textbf{方向}，不能替代判别式或 $S$ 给出的\textbf{真实存在性}。
\par\medskip\noindent
\textbf{考场 10 秒路径：}
\[
\boxed{\text{有中点}\;\Longrightarrow\;\text{两点平方相减}\;\Longrightarrow\;
\frac{x_0\Delta x}{a^2}=\frac{y_0\Delta y}{b^2}}.
\]
当 $y_0\ne0$ 且真实存在非退化弦时，
\[
\boxed{k=\frac{b^2x_0}{a^2y_0}}.
\]
\textbf{反向求弦：}设 $S=x_0^2/a^2-y_0^2/b^2$，$M\ne O$ 时
\[
\boxed{\text{候选弦}\ \frac{x_0x}{a^2}-\frac{y_0y}{b^2}=S},\qquad
\boxed{\text{真有弦}\iff S<0\ \text{或}\ S>1}.
\]
$S<0$ 为异支；$S>1$ 为同支。$S=1$ 只有切点；$0\le S<1$ 一般无弦，\textbf{但原点 $O$ 例外}：中心弦斜率满足 $|k|<b/a$。
\par\medskip\noindent
\textbf{一句话抓本质：}相减求\textbf{方向}，相加求\textbf{半弦位移的平方关系}；
\textbf{方向要对，沿这个方向的实距离也要存在}。$S$ 判\textbf{存在}；原点和竖直弦要\textbf{另看}。
\end{summarybox}
